% =============================================================================
% Capítulo 4: Estadísticas de Procesamiento (nombre literal del enunciado)
% Avance N°2. Las tablas las genera herramientas/estadisticas_etl.py desde el
% registro de Apache Hop y la base, después de la ejecución final del ETL.
% =============================================================================
\section{Estadísticas de Procesamiento}
\label{sec:estadisticas}

Este capítulo presenta los resultados del proceso ETL que puebla el
Datawarehouse: cuántos registros se leyeron de cada fuente, cuántos llegaron
al staging, cuántos se rechazaron y por qué, cuántos se cargaron en el modelo
estrella y cuánto tardó cada etapa. Las cifras corresponden a la ejecución
final del ETL, hecha con el comando \bd{python etl/ejecutar_etl.py
--preparar} sobre PostgreSQL 17.11 y Apache Hop 2.20.0. Las tablas de este
capítulo no se escriben a mano: las genera el script
\bd{herramientas/estadisticas_etl.py} desde el registro de ejecución de Hop y
desde la propia base.

\subsection{Ejecución del Proceso}

El ETL se ejecuta con un solo comando, que corre la preparación en Python y
luego el workflow principal de Apache Hop
(Sección~\ref{sec:flujo_etl}). La Figura~\ref{fig:hop_metricas} muestra las
métricas de una ejecución en Apache Hop GUI: filas leídas, escritas y con
error de cada transformación del pipeline de hechos. La
Figura~\ref{fig:etl_terminal} muestra el final de la ejecución por línea de
comandos, con el resultado de cada etapa del workflow.

\captura{07d_hop_metricas_ejecucion.png}{Métricas de ejecución del pipeline \bd{03_hechos} en Apache Hop GUI\label{fig:hop_metricas}}
\captura{07e_terminal_ejecucion_etl.png}{Ejecución del ETL completo por línea de comandos (\bd{python etl/ejecutar_etl.py})\label{fig:etl_terminal}}

\subsection{Registros por Etapa y Línea de Negocio}

La Tabla~\ref{tab:est_lineas} resume el recorrido de los registros. De las
fuentes se leyeron más de 80 millones de filas, casi todas de la encuesta de
pasajes aéreos de BTS, que se lee completa para muestrear sus diez trimestres.
De ellas se preparó una muestra de 120.000 transacciones, que pasó al staging
sin pérdidas. El 98,6~\% se cargó en \bd{fact_venta} y el resto se rechazó
por no cumplir alguna regla de calidad.

% Generado por herramientas/estadisticas_etl.py. No editar a mano.
\begin{table}[!htbp]
  \centering
  \caption{Registros procesados por línea de negocio}
  \label{tab:est_lineas}
  \small
  \begin{tabularx}{\linewidth}{@{}>{\raggedright\arraybackslash}p{2.6cm}>{\raggedright\arraybackslash}Xrrrrr@{}}
    \toprule
    \textbf{Línea} & \textbf{Fuente} & \textbf{Leídos} & \textbf{Staging} & \textbf{Rechazados} & \textbf{Cargados} & \textbf{\% carga} \\
    \midrule
    Buses InterUrbanos & Supertransporte & 1.260.343 & 25.000 & 0 & 25.000 & 100,0 \% \\
    Trenes & UK Train Rides & 31.653 & 15.000 & 0 & 15.000 & 100,0 \% \\
    Radio Taxis & NYC TLC & 2.038.653 & 25.000 & 273 & 24.727 & 98,9 \% \\
    Logística y Distribución & Olist & 99.441 & 20.000 & 199 & 19.801 & 99,0 \% \\
    Aerolíneas & BTS DB1B & 77.949.326 & 15.000 & 303 & 14.697 & 98,0 \% \\
    URapid Delivery & Delivery Center & 368.999 & 20.000 & 962 & 19.038 & 95,2 \% \\
    \midrule
    \textbf{Total} &  & \textbf{81.748.415} & \textbf{120.000} & \textbf{1.737} & \textbf{118.263} & \textbf{98,6 \%} \\
    \bottomrule
  \end{tabularx}
  \par\smallskip\footnotesize Leídos: filas de la fuente leídas por la preparación en Python. Staging: muestra cargada en stg.transaccion. Cargados: filas de fact\_venta. El porcentaje de carga se calcula sobre staging.
\end{table}


Buses y trenes no tienen rechazos. En buses, la preparación ya descarta lo que
no es una venta: despachos sin pasajeros, vehículos que no son buses (como
automóviles o camionetas) y rutas con municipios inexistentes. En trenes, el
dataset es ficticio y no trae registros inválidos. Delivery tiene la tasa de
carga más baja (95,2~\%), porque el 4,7~\% de sus pedidos aparece cancelado
en la fuente.

\subsection{Rechazos y Calidad de Datos}

Cada transacción pasa por las reglas de calidad de la vista
\bd{stg.v_transaccion_conformada} (Figura~\ref{fig:vista_reglas}). La primera
regla que no se cumple queda como motivo de rechazo, y la transacción se
registra en \bd{stg.rechazo} en vez de cargarse. La
Tabla~\ref{tab:est_rechazos} muestra los motivos encontrados.

% Generado por herramientas/estadisticas_etl.py. No editar a mano.
\begin{table}[!htbp]
  \centering
  \caption{Transacciones rechazadas por motivo}
  \label{tab:est_rechazos}
  \small
  \begin{tabularx}{\linewidth}{@{}>{\raggedright\arraybackslash}Xlr@{}}
    \toprule
    \textbf{Motivo (regla de calidad)} & \textbf{Línea} & \textbf{Filas} \\
    \midrule
    estado en la fuente: cancelado & URapid Delivery & 943 \\
    cantidad fuera de 1 a 9 & Aerolíneas & 293 \\
    estado en la fuente: sin\_cobro & Radio Taxis & 166 \\
    estado en la fuente: cancelado & Logística y Distribución & 106 \\
    monto bruto no positivo & Logística y Distribución & 90 \\
    estado en la fuente: disputa & Radio Taxis & 48 \\
    monto bruto no positivo & Radio Taxis & 25 \\
    duración negativa & Radio Taxis & 21 \\
    distancia atípica & URapid Delivery & 18 \\
    distancia atípica & Radio Taxis & 12 \\
    monto bruto no positivo & Aerolíneas & 8 \\
    cantidad fuera de 1 a 9 & Logística y Distribución & 3 \\
    estado en la fuente: tarifa\_masiva & Aerolíneas & 2 \\
    monto atípico & Radio Taxis & 1 \\
    monto atípico & URapid Delivery & 1 \\
    \midrule
    \textbf{Total} &  & \textbf{1.737} \\
    \bottomrule
  \end{tabularx}
\end{table}


Los rechazos corresponden a problemas reales de las fuentes:

\begin{itemize}[leftmargin=*]
  \item \textbf{Estado en la fuente:} pedidos cancelados (delivery y
        logística), viajes de taxi sin cobro o con el pago en disputa, y
        pasajes aéreos de tarifa masiva. No son ventas.
  \item \textbf{Cantidad fuera de 1 a 9:} registros de BTS que agrupan a
        decenas de pasajeros en un mismo itinerario, propios de vuelos chárter
        o grupos, que distorsionarían el ticket promedio.
  \item \textbf{Monto bruto no positivo:} pasajes aéreos con tarifa 0, envíos
        sin flete y viajes de taxi con monto cero o negativo.
  \item \textbf{Duración negativa, distancia o monto atípicos:} viajes de taxi
        que terminan antes de empezar, recorridos de cientos de kilómetros
        dentro de Nueva York, entregas de delivery a más de 50 km o pedidos de
        más de un millón de pesos.
\end{itemize}

Ninguna transacción se rechazó por un código sin equivalencia, un cliente
inexistente o una moneda sin tipo de cambio: los catálogos de
\bd{etl/referencias/} cubren todos los valores de las fuentes.

\subsection{Detalle por Transformación}

La Tabla~\ref{tab:est_hop} reproduce las filas que informó Apache Hop para cada
transformación de los tres pipelines. En el pipeline de hechos se ve el
reparto del filtro de calidad: las filas que no cumplen las reglas van a
\bd{Registrar rechazo}, y las demás pasan por las seis búsquedas de llaves
sustitutas hasta \bd{Cargar fact_venta}. Ninguna transformación registró
errores.

% Generado por herramientas/estadisticas_etl.py. No editar a mano.
\begin{table}[H]
  \centering
  \caption{Filas por transformación en la ejecución del ETL (registro de Apache Hop)}
  \label{tab:est_hop}
  \small
  \begin{tabularx}{\linewidth}{@{}l>{\raggedright\arraybackslash}Xrrr@{}}
    \toprule
    \textbf{Pipeline} & \textbf{Transformación} & \textbf{Entrada} & \textbf{Salida} & \textbf{Errores} \\
    \midrule
    01\_staging & Leer tipo\_cambio.csv & 13 & 12 & 0 \\
     & Leer equivalencia.csv & 50 & 49 & 0 \\
     & Cargar stg.tipo\_cambio & 12 & 12 & 0 \\
     & Cargar stg.equivalencia & 49 & 49 & 0 \\
     & Leer zonas.csv & 3.654 & 3.653 & 0 \\
     & Cargar stg.zona & 3.653 & 3.653 & 0 \\
     & Leer clientes.csv & 10.001 & 10.000 & 0 \\
     & Cargar stg.cliente & 10.000 & 10.000 & 0 \\
     & Leer transacciones.csv & 120.001 & 120.000 & 0 \\
     & Cargar stg.transaccion & 120.000 & 120.000 & 0 \\
    \midrule
    02\_dimensiones & Leer canal\_venta.csv & 8 & 7 & 0 \\
     & Leer servicio.csv & 18 & 17 & 0 \\
     & Leer medio\_pago.csv & 11 & 10 & 0 \\
     & Cargar dim\_canal\_venta & 7 & 7 & 0 \\
     & Cargar dim\_medio\_pago & 10 & 10 & 0 \\
     & Cargar dim\_servicio & 17 & 17 & 0 \\
     & Leer stg.zona & 3.653 & 3.653 & 0 \\
     & Leer stg.cliente & 10.000 & 10.000 & 0 \\
     & Cargar dim\_zona & 3.653 & 3.653 & 0 \\
     & Cargar dim\_cliente & 10.000 & 10.000 & 0 \\
    \midrule
    03\_hechos & Leer transacciones conformadas & 120.000 & 120.000 & 0 \\
     & ¿Cumple las reglas? & 120.000 & 120.000 & 0 \\
     & Registrar rechazo & 1.737 & 1.737 & 0 \\
     & Buscar id\_cliente & 118.263 & 118.263 & 0 \\
     & Buscar id\_servicio & 118.263 & 118.263 & 0 \\
     & Buscar id\_zona\_origen & 118.263 & 118.263 & 0 \\
     & Buscar id\_zona\_destino & 118.263 & 118.263 & 0 \\
     & Buscar id\_canal\_venta & 118.263 & 118.263 & 0 \\
     & Buscar id\_medio\_pago & 118.263 & 118.263 & 0 \\
     & Cargar fact\_venta & 118.263 & 118.263 & 0 \\
    \bottomrule
  \end{tabularx}
  \par\smallskip\footnotesize Entrada: filas leídas del archivo, de la base o de la transformación anterior. Salida: filas escritas en la base o enviadas a la transformación siguiente. En los archivos CSV, la entrada incluye la fila de encabezado.
\end{table}


\subsection{Tiempos de Ejecución}

La Tabla~\ref{tab:est_tiempos} muestra la duración de cada etapa. La
preparación en Python está dominada por la lectura de los archivos de BTS
(cerca de 78 millones de filas). El workflow de Hop completo, desde la creación
del staging hasta la validación final, tarda menos de 20 segundos.

% Generado con herramientas/estadisticas_etl.py; tiempos de Hop fijados a la ejecución de la Figura 4.2 (18-11-2026 21:26).
\begin{table}[!htbp]
  \centering
  \caption{Tiempo de ejecución por etapa}
  \label{tab:est_tiempos}
  \small
  \begin{tabularx}{\linewidth}{@{}>{\raggedright\arraybackslash}Xr@{}}
    \toprule
    \textbf{Etapa} & \textbf{Segundos} \\
    \midrule
    Preparación en Python: Buses InterUrbanos & 4,4 \\
    Preparación en Python: Trenes & 0,1 \\
    Preparación en Python: Radio Taxis & 0,4 \\
    Preparación en Python: Logística y Distribución & 1,4 \\
    Preparación en Python: Aerolíneas & 90,6 \\
    Preparación en Python: URapid Delivery & 0,9 \\
    \midrule
    \textit{Subtotal preparación en Python} & \textit{97,8} \\
    \midrule
    Hop: pipeline 01\_staging & 3,1 \\
    Hop: pipeline 02\_dimensiones & 0,3 \\
    Hop: pipeline 03\_hechos & 10,6 \\
    \midrule
    \textit{Subtotal workflow principal de Hop} & \textit{15,0} \\
    \midrule
    \textbf{Total del ETL} & \textbf{112,8} \\
    \bottomrule
  \end{tabularx}
  \par\smallskip\footnotesize El subtotal de Hop es la duración del workflow principal: además de los tres pipelines, ejecuta los scripts SQL de staging, de preparación de la carga y de validación. El total suma la preparación de las seis líneas y el workflow. Los tiempos de Python provienen del registro de la preparación; los de Hop, de la ejecución de la Figura~\ref{fig:etl_terminal}.
\end{table}


\subsection{Resultado de la Carga}

La Tabla~\ref{tab:est_tablas} muestra las filas de cada tabla al terminar la
carga, y las figuras siguientes, las consultas de verificación de
\bd{sql/98_consultas_carga.sql} ejecutadas en pgAdmin.

% Generado por herramientas/estadisticas_etl.py. No editar a mano.
\begin{table}[H]
  \centering
  \caption{Filas por tabla después de la carga}
  \label{tab:est_tablas}
  \small
  \begin{tabularx}{\linewidth}{@{}>{\raggedright\arraybackslash}Xr@{}}
    \toprule
    \textbf{Tabla} & \textbf{Filas} \\
    \midrule
    \bd{stg.cliente} & 10.000 \\
    \bd{stg.equivalencia} & 49 \\
    \bd{stg.rechazo} & 1.737 \\
    \bd{stg.tipo\_cambio} & 12 \\
    \bd{stg.transaccion} & 120.000 \\
    \bd{stg.zona} & 3.653 \\
    \bd{dwh.dim\_canal\_venta} & 8 \\
    \bd{dwh.dim\_cliente} & 10.000 \\
    \bd{dwh.dim\_medio\_pago} & 11 \\
    \bd{dwh.dim\_servicio} & 17 \\
    \bd{dwh.dim\_tiempo} & 1.096 \\
    \bd{dwh.dim\_zona} & 3.653 \\
    \bd{dwh.fact\_venta} & 118.263 \\
    \bottomrule
  \end{tabularx}
\end{table}


La consistencia de la carga se comprobó de tres formas:

\begin{enumerate}[leftmargin=*]
  \item \textbf{Validación dentro del ETL:} la última acción del workflow
        (\bd{sql/04_validar_carga.sql}) comprueba que cada transacción del
        staging terminó en \bd{fact_venta} o en \bd{stg.rechazo}, y que por
        cada línea la cantidad y el monto neto de las transacciones válidas
        del staging son iguales a los de \bd{fact_venta}. Si algo no cuadra,
        el ETL termina con error.
  \item \textbf{Conciliación en pgAdmin:} la Figura~\ref{fig:carga_conciliacion}
        muestra las diferencias de filas y de monto neto entre el staging y
        \bd{fact_venta} por línea de negocio: todas son cero.
  \item \textbf{Pruebas automáticas:} las 10 pruebas de
        \bd{tests/test_carga.py} verifican, entre otras cosas, que
        \bd{dim_tiempo} tiene los 1.096 días de 2023 a 2025 sin huecos, que la
        edad de cada venta coincide con la fecha de nacimiento del cliente, que
        ninguna venta es anterior a la fecha de registro de su cliente y que
        todos los clientes están marcados como sintéticos. Con las 29 pruebas
        del DDL y las 4 de la preparación, se aprobaron 43 de 43.
\end{enumerate}

\captura{06a_carga_filas_por_tabla.png}{Filas por tabla después de la carga (consulta 1 de \bd{sql/98_consultas_carga.sql})\label{fig:carga_filas}}
\captura{06b_carga_por_linea.png}{Transacciones en staging, rechazadas y cargadas por línea (consulta 2)\label{fig:carga_lineas}}
\captura{06c_carga_rechazos.png}{Transacciones rechazadas por motivo (consulta 3)\label{fig:carga_rechazos}}
\captura{06d_carga_conciliacion.png}{Conciliación por línea entre el staging y \bd{fact_venta}: diferencias en cero (consulta 4)\label{fig:carga_conciliacion}}

\subsection{Primera Lectura de los Datos Cargados}

Con el Datawarehouse poblado, las consultas 5 y 6 de
\bd{sql/98_consultas_carga.sql} entregan un adelanto de los indicadores que el
cubo y el dashboard desarrollarán en el Avance N°3
(Figuras~\ref{fig:carga_ventas} y~\ref{fig:carga_indicadores}). Las ventas por
línea y año permiten comparar el aporte de cada línea, y los indicadores por
año calculan los KPI definidos en el modelo: clientes activos, frecuencia de
compra, gasto anual por cliente, tasa de descuento y clientes multilínea.

Estas cifras deben leerse con las limitaciones de la
Sección~\ref{sec:limitaciones}: los montos provienen de fuentes de distintos
países convertidas a pesos chilenos, y el comportamiento de los clientes entre
líneas proviene del generador de clientes sintéticos. Por ejemplo, la
proporción de clientes multilínea depende de los supuestos de afinidad del
generador, no de clientes reales.

\captura{06e_carga_ventas_linea_anio.png}{Ventas por línea de negocio y año (consulta 5)\label{fig:carga_ventas}}
\captura{06f_carga_indicadores_anio.png}{Indicadores por año: clientes activos, frecuencia, gasto, descuento y clientes multilínea (consulta 6)\label{fig:carga_indicadores}}
